\section{导入 \projectname.cls文档类及相关配置宏包}
\subsection{导入\projectname.cls文档类}
\begin{lstlisting}[language=tex]
    \documentclass[young,application,twoside]{style/nsfc}
    %- Multiple optional arguments:
    %- [<application|closeout>]: 申请报告与结题报告
    %- [<young|general>]: 青年项目与面上项目
    %- [<oneside|twoside|null>]% oneside eprint, twoside eprint, or paper print
    %- [fontset=<adobe|none|...>]% specify font set instead of automatic detection
    %- [scheme=plain]% thesis writing of international students
    %- [draft]% show draft version information
    %- [standard options for ctex book class: draft|paper size|font size|...]%
    %---------------------------------------------------------------------------%
 \end{lstlisting}%
在导入 \projectname.cls文档类时，提供了下述选项。
\begin{enumerate}
    \item \emph{指定文档类型}：在引入文档类时，可以从两选项 “[<application|closeout>]” 中选一个来指定文档是申请书或结题报告（ \emph{结题报告待开发}），即“申请书application与“结题报告closeout”。进一步地，需要通过选项 “[<young|general>]” 指定青年项目或者是面上项目。这两组选项配合可完全确定文档类型。
    \item \emph{打印模式}: 在引入文档类时，可以从三选项“[<oneside|twoside|null>]”中选一个来指定文档的打印模式，null为默认的缺省选项，不需要填写。oneside和twoside依次对应单面打印和双面打印。oneside和twoside均没有给出时，采用等边距空白。
    \item \emph{平台字体族选项}：模板默认自动检测操作系统平台，根据操作系统平台选择字体族，在引入文档类时，不需要指定“fontset”选项。当然，我们也可以直接通过明确指定选项“[<fontset=<windows|mac|adobe|none|>]”强制特定操作系统平台类型，以达到手动确定操作系统平台字体族。详见章节\ref{sec:fontset}。
    \item \emph{启用手稿模式}：在引入文档类时，填入选项“[<draft>]”，模板会启用draft模式编译文档，并在文档第二页打印“Draft Version(日期)”。
    \item \emph{其它选项}: 其它受ctex book class支持的选项，包括: "[<draft|paper size|font size|...>]"。
\end{enumerate}



\subsection{导入基础配置宏包cnbasesetup.sty}\label{ssec:cnbasesetup}
\begin{lstlisting}[language=tex]
    \usepackage[super,list,table]{style/cnbasesetup}% document settings
    %- usage: \usepackage[option1,option2,...,optionN]{style/cnbasesetup}
    %- Multiple optional arguments:
    %- [bibauto|bibtex|biber]% set bibliography processor and package
    %- [<numbers|super|authoryear|alpha>]% set citation and reference style
    %   - <numbers>: textual: Jones [1]; parenthetical: [1]
    %   - <super>: textual: Jones superscript [1]; parenthetical: superscript [1]
    %   - <authoryear>: textual: Jones (1995); parenthetical: (Jones, 1995)
    %
    %- [lscape]% provide landscape layout environment
    %- [color]% provide color support via xcolor package
    %- [tikz]% provide complex diagrams via tikz package
    %- [table]% provide complex tables via ctable package
    %- [list]% provide enhanced list environments for algorithm and coding
    %- [math]% enable some extra math packages
    %- [xlink]% disable link colors
\end{lstlisting}%
\begin{enumerate}
    \item \emph{指定参考文献处理器}：在引入cnbasesetup宏包时，我们可以直接通过指定三个选项"[bibauto|bibtex|biber]"中的任意一个来确定学位论文参考文献处理器。
        \begin{enumerate}
            %\footnotesize
            \item "bibtex"选项指定参考文献处理器为bibtex。当参考文献处理器选择为“bibtex”时，在需要的位置直接使用带参考文献库名参数的命令 \verb|\nsfcbibliography{<bibname>}| 打印参考文献列表，这里没有文献库扩展名“.bib”。
            \item “biber”选项指定参考文献处理器为“biblatex”。当参考文献处理器选择为“biber”时，需要在导言区添加命令 \verb|\addbibresource{<bibname.bib>}| 以指定参考文献库为"bibname.bib"，并且这个命令需要带上扩展名“.bib”，然后在需要的位置直接使用不带参数的命令“\backslash \projectname bibliography”打印参考文献列表。
            \item “bibauto”为文档自定义的选项，参考文献处理器由 \projectname.cls文档类指定。这也是缺省默认的处理器模式， \projectname.cls默认配置为biblatex。 \emph{该模式下，不可自定义参考文献库文件，仅可采用默认参考文献库文件“ref.bib”}。
        \end{enumerate}

    没有明确指定参考文献库名的情况下，模板默认配置参考文献库名为“ref.bib”。样例文档中文献处理器没有采用默认配置的“bibauto”，而是设置为“biber”，并且通过命令 \verb|\addbibresource{nsfc.bib}| 自定义了参考文献库文件为“nsfc.bib”。
    
    对于需要打印多段参考文献列表的，只能选择“biber”或“bibauto”作为参考文献处理器，比如引入“综述”及“综述参考文献”。
    %
    \item \emph{配置参考文献引用风格}: 在引入cnbasesetup宏包时，我们可以直接通过指定三个选项“[<numbers|super|authoryear>]”中的任意一个来确定参考文献引用风格。
        \begin{enumerate}
        %\footnotesize
        \item \verb|number|：按出现顺序排序，以数字方式索引引用；
        \item \verb|super|：按出现顺序排序，以数字方式索引引用，并且引用处采用上角标小数字引用；
        \item \verb|authoryear|： 以“作者 年份”方式索引引用；
        \item \verb|nsfc|：指定国家自然基金委模板中提到的引用格式，以小括号格式标注文献顺序。
        \end{enumerate}
    \item \emph{局部转置版面宏包}： 参数"[<lscap>]"指定文档引入宏包fancyhdr、pdflscape、textpos，提供了landscape环境。
    \item \emph{颜色宏包}： 参数"[<color>]"指定文档引入宏包xcolor，提供丰富的颜色支持。
    \item \emph{tikz绘图宏包}： 参数"[<tikz>]"指定文档引入额外扩展的绘图宏包tikz ，提供绘图支持。
    \item \emph{table类宏包}： 参数"[<table>]"指定文档引入宏包ctable、multirow、longtable、makecell，提供绘图支持。
    \item \emph{list类宏包}： 参数"[<list>]"指定文档引入宏包:
        \begin{lstlisting}[language=tex]
            \RequirePackage{verbatim}% improve verbatim environment
            \RequirePackage{enumitem}% configure the enumerate environment
            \setlist[enumerate]{wide=\parindent}% only indent the first line
            \setlist[itemize]{wide=\parindent}% only indent the first line
            \setlist{nosep}% default text spacing
            \RequirePackage{listings}% source code
            \RequirePackage{algpseudocode,algorithm,algorithmicx}% algorithm
            \providecommand{\algname}{Algorithm}%
            \def\ALG@name{\algname}% rename label
        \end{lstlisting}%
    \item \emph{math宏包}： 参数"[<math>]"指定文档引入额外数学宏包“mathtools”。
    \item \emph{禁用链接颜色}： 参数"[<xlink>]"指定文档禁用链接的颜色。
\end{enumerate}


\subsection{导入模板配置宏包nsfcsetup.sty}
\begin{lstlisting}[language=tex]
        \usepackage{style/nsfcsetup}% document settings
 \end{lstlisting}%



 
 \section{\projectname 学位论文信息配置命令}
 建议把学位论文信息配置命令写入文件“\projectname /Frontinfo.tex”文件内。下面详细列出各个配置命令。
 \begin{enumerate}
    \item \emph{\backslash nsfclogo\{\}}：配置NSFC logo。例如：\newline \verb|\nsfclogo[width=4.87cm,height=2.41cm]{official/nsfc/nsfc-logo}|。
    \item \emph{\backslash category\{\}}：资助类别：一般面上或青年基金; 一般不用填，除非默认资助类别显示的不是你所期望的，例如：\verb|\category{一般项目}|。
    \item \emph{\backslash subcategory\{\}}：亚类说明，例如： \verb|\subcategory{亚类说明测试}|。
    \item \emph{\backslash additionalnotes\{\}}：附注说明，例如： \verb|\additionalnotes{附注说明测试}|。
    \item \emph{\backslash title\{\}}：项目中文名称，例如： \verb|\title{脉冲神经元模型}|。可使用 \verb|\\| 命令手动控制换行，如：\verb|\title{脉冲神经网\\高效学习理论研究}|。
    \item \emph{\backslash title*\{\}}：项目英文名称，例如： \verb|\title*{Spiking Neuron Model}|。
    \item \emph{\backslash author\{\}}：申请人，例如：\verb|\author{朱昌波}|。
    \item \emph{\backslash email\{\}}：申请人邮件，例如： \verb|\email{1207146285@qq.com}|。
    \item \emph{\backslash postcode\{\}}：申请人邮编，例如： \verb|\postcode{233222}|。
    \item \emph{\backslash phone\{\}}：申请人电话，例如： \verb|\phone{12345678901}|。
    \item \emph{\backslash address\{\}}：申请人地址，例如： \verb|\address{某某省某某市某某街19号某某大学某某学院}|。
    \item \emph{\backslash institute\{\}}：申请人单位，例如： \verb|\institute{某某大学}|。
    \item \emph{\backslash institutephone*\{\}}：申请人单位电话，例如： \verb|\institutephone{0230-23456789}|。
    \item \emph{\backslash researchfield\{\}}：本申请研究领域或方向，例如： \verb|\researchfield{机器学习、计算神经科学}|。
    \item \emph{\backslash degree\{\}}：申请人最高学位，例如： \verb|\degree{工学博士}|。
\end{enumerate}